\documentclass[11pt]{scrartcl}

\usepackage[stixtwo]{fontsetup}
\usepackage{amsmath}

\begin{document}
A köbgyökvonás megtartja az előjelet, a $-4$-gyel szorzás megfordítja, így
$x_n$ előjele lépésenként váltakozik: $x_0 = -1 < 0$, $x_1 = 4 > 0$, és így
tovább. Az abszolút értékre $|x_n| = 4 \sqrt[3]{|x_{n - 1}|}$.

Legyen $y_n = \log_4 |x_n|$. Ekkor $y_0 = 0$ és
$y_n = 1 + \frac{1}{3} y_{n - 1}$, azaz
\(
y_n - \frac{3}{2} = \frac{1}{3} \left( y_{n - 1} - \frac{3}{2} \right)
\),
így
\[
	y_n = \frac{3}{2} \left( 1 - \frac{1}{3^n} \right) \to \frac{3}{2},
	\qquad
	|x_n| = 4^{y_n} \to 4^{\frac{3}{2}} = 8.
\]
Páratlan $n$-re $x_n = |x_n| \to 8$, páros $n$-re $x_n = -|x_n| \to -8$.

A sorozat torlódási pontjai tehát: $\{-8, 8\}$.
\end{document}
